\SourceLeaf{298}
Aus $\dfrac{6^\ast\times 8^\ast}{11^\ast}$ oder $\dfrac{6^\ast\times 10^\ast}{9^\ast}$ oder $\dfrac{8^\ast\times 10^\ast}{7^\ast}$ folgt

\[
PP=qq,
\]

woraus nothwendig

\[
P=+q
\]

geschlossen werden muss, in so fern man nur Dreiecke betrachtet, wo Seiten und Winkel nicht über $180^0$ hinausgehen. Dann hat man ohne weiteres

\[
Q=+s,\qquad R=+p,\qquad S=+r,
\]

welches zusammen die vier gewünschten Sätze sind.

Ich überlasse es Ihnen, in wie fern Sie hiervon, oder von andern mir gewöhnlichen Rechnungsanordnungen, in Ihrem Werke Gebrauch machen wollen. \mbox{.\,.\,.}

\medskip

\Needspace{8\baselineskip}
\begin{center}
[2.]

\textsc{Gauss} an \textsc{Gerling}. Göttingen, 26. Juni 1816.
\end{center}

\mbox{.\,.\,.\,.\,.\,.\,.} Ich habe bisher immer so vielfache Zerstreuungen gehabt, dass ich noch nicht habe dazu kommen können, über die zweckmässigste Ordnung der ersten Sätze der sphärischen Trigonometrie, bei einer geometrischen Behandlung, nachzudenken. Ich sollte indessen glauben, dass es am einfachsten sein würde, Euklid XI. 20. 21 zum Grunde zu legen und den Satz, dass $A+B+C>180^0$ ist, aus dem Polardreieck abzuleiten. Ich glaube kaum, dass sich aus der Betrachtung des ebenen Dreiecks zwischen den Winkelpunkten des sphärischen irgend ein erheblicher Gewinn ziehen lässt, da es keine \textls{sehr} einfache Relationen zwischen dessen Winkeln zu geben scheint. Folgende möchte wohl mit die zierlichste sein: wenn $A'$, $B'$, $C'$ die Winkel des ebenen Dreiecks sind, und

\[
A+B+C-180^0=u,
\]

so ist

\[
\cos A'=\cos\tfrac{1}{2}a\cos(A-\tfrac{1}{2}u).
\]

\mbox{.\,.\,.\,.\,.\,.\,.\,.\,.}
